\label{ch:vacuum}

\section{Derivaci\'on de la escala de acci\'on a partir del estado dodecaf\'asico}

La d\'ecada de acci\'on se determina intr\'insecamente, sin adoptarla como
unidad heredada. El funcional determinantal local se aplica a la salida HMT
previamente construida \(\alpha_{\HMT}\) y define
\begin{equation}
\Sigma_H=\varphi_{\HMT}\alpha_{\HMT}^{16}.
\label{eq:action-determinantal-reader}
\end{equation}
El exponente \(16\) es el orden del cokernel local del cierre firmado; su valoraci\'on decimal fija
\begin{equation}
\operatorname{ord}_{10}(\Sigma_H)=-34.
\label{eq:action-decade}
\end{equation}

El origen del exponente se puede verificar dentro de este volumen. Sobre la
\'orbita local de cuatro hojas act\'ua el bloque de Hadamard
\begin{equation}
H_4=
\begin{pmatrix}
1&1&1&1\\
1&1&-1&-1\\
1&-1&1&-1\\
1&-1&-1&1
\end{pmatrix}.
\label{eq:action-hadamard}
\end{equation}

\begin{lemma}[\'Indice local de la recta de acci\'on]
La forma normal de Smith de \(H_4\) es
\begin{equation}
\operatorname{SNF}(H_4)=\operatorname{diag}(1,2,2,4).
\label{eq:action-smith}
\end{equation}
En consecuencia,
\[
\operatorname{coker}(H_4)\simeq
\Z/2\Z\oplus\Z/2\Z\oplus\Z/4\Z,
\qquad
|\operatorname{coker}(H_4)|=16.
\]
\end{lemma}

\begin{proof}
El m\'aximo com\'un divisor de las entradas es uno; el de los menores
\(2\times2\) no nulos es dos; el de los menores \(3\times3\) es cuatro; y
\(|\det H_4|=16\). Los cocientes sucesivos de estos divisores determinantes
son \(1,2,2,4\), que prueban \cref{eq:action-smith}. El producto de los
invariantes da el orden del con\'ucleo.
\end{proof}

La norma de la secci\'on \(\alpha_{\HMT}\) sobre esas diecis\'eis hojas
produce \(\alpha_{\HMT}^{16}\); la autoescala \(\varphi_{\HMT}\) act\'ua
una sola vez sobre la base, no una vez por hoja. De ah\'i
\cref{eq:action-determinantal-reader}. Las comparaciones certificadas
\begin{equation}
\alpha_{\HMT}^{16}<10^{-34}
<\varphi_{\HMT}\alpha_{\HMT}^{16}<10^{-33}
\label{eq:action-decade-bounds}
\end{equation}
demuestran, sin introducir la unidad SI, que el orden decimal es
exactamente \(-34\). La evaluaci\'on terminal comienza por
\begin{equation}
\Sigma_H
=1.046243838104756580184229208303089\ldots\times10^{-34}.
\label{eq:action-sigma-value}
\end{equation}
La sustituci\'on del con\'ucleo por tres copias con exponente \(16^3\), o la
aplicaci\'on reiterada de \(\varphi\) en cada hoja, modificar\'ia la
genealog\'ia. Ambas operaciones constituyen criterios de refutaci\'on y no
convenciones equivalentes.
La cadena causal que se conserva es, por tanto,
\[
\APP\to\TRIT\to\TPK\to U_{12}^{\rm sign}
\to\alpha_{\HMT}\to\Sigma_H
\to10^{-34}\U_S\to
\{\hbar_{\rm pre},\hbar_{\rm ret}\}.
\]
Este volumen no rederiva \(\alpha_{\HMT}\), sino que desarrolla la geometr\'ia
metrol\'ogica posterior a la selecci\'on del tipo y de la d\'ecada por
\(\Sigma_H\). El entero \(-34\) no es una mantisa ni un ajuste al Sistema
Internacional, sino la valoraci\'on terminal del funcional local.

\section{Secciones orientadas de la recta de acci\'on}

Las realizaciones anterior y posterior al retorno constituyen dos secciones
de una misma recta de acci\'on:

\begin{equation}
\hbar_{\rm pre}=H_5\,10^{-34}\U_S,\qquad
\hbar_{\rm ret}=\eta_{\rm ret}\,10^{-34}\U_S.
\label{eq:action-twoway}
\end{equation}

La d\'ecada com\'un determina la recta, mientras que la informaci\'on
orientada se expresa mediante

\begin{equation}
R_{\rm act}=\frac{\hbar_{\rm pre}}{\hbar_{\rm ret}}
=\frac{H_5}{\eta_{\rm ret}},
\qquad
\xi_{\rm act}=\frac12\log R_{\rm act}.
\label{eq:action-rapidity}
\end{equation}

La primera cantidad proporciona una coordenada multiplicativa y la segunda
una coordenada aditiva de tipo rapidez. No representan dos constantes de
Planck independientes, sino las dos orientaciones de una misma secci\'on
respecto del retorno.

Si \(h_a=2\pi\hbar_a\), \(a\in\{\mathrm{pre},\mathrm{ret}\}\), una magnitud proporcional a \(h_a^{1/2}\) cambia como \(R_{\rm act}^{1/2}\), una proporcional a \(h_a^{-1/2}\) cambia como \(R_{\rm act}^{-1/2}\), y un cociente donde la acci\'on se cancela es bidireccional invariante.

